\import{topology/topological-ring-boundedness.tex}
\import{algebra/monoid/subset.tex}
\import{algebra/ring/binomial.tex}
\import{algebra/group/finite-sum.tex}

\section{Power-bounded subsets and elements}

\begin{definition}\label{power_bounded_set}
    $\PowerBounded{R}{B}$ iff
        $R$ is a commutative topological ring and
        $B\subseteq\carrier[R]$ and
        $\Bounded{R}{\pointwiseSubsetPowers{
            \MultiplicativeMonoid{R}}{B}}$.
\end{definition}

\begin{proposition}\label{power_bounded_set_iff}
    Let $R$ be a commutative topological ring.
    Then $\PowerBounded{R}{B}$ iff
        $B\subseteq\carrier[R]$ and
        $\Bounded{R}{\pointwiseSubsetPowers{
            \MultiplicativeMonoid{R}}{B}}$.
\end{proposition}
\begin{proof}
    Follows by \auto{power_bounded_set}.
\end{proof}

\begin{proposition}\label{power_bounded_set_carrier}
    Suppose $\PowerBounded{R}{B}$.
    Then $B\subseteq\carrier[R]$.
\end{proposition}
\begin{proof}
    Follows by \obv{power_bounded_set}.
\end{proof}

\begin{proposition}\label{power_bounded_set_powers_bounded}
    Suppose $\PowerBounded{R}{B}$.
    Then $\Bounded{R}{\pointwiseSubsetPowers{
        \MultiplicativeMonoid{R}}{B}}$.
\end{proposition}
\begin{proof}
    Follows by \obv{power_bounded_set}.
\end{proof}

\begin{theorem}\label{power_bounded_set_downward}
    Suppose $\PowerBounded{R}{B}$ and $A\subseteq B$.
    Then $\PowerBounded{R}{A}$.
\end{theorem}
\begin{proof}
    Hence $R$ is a commutative ring by
        \obv{commutative_topological_ring_is_commutative_ring}.
    Hence $R$ is a ring by \obv{commutative_ring_is_ring}.
    Let $M=\MultiplicativeMonoid{R}$.
    $M$ is a monoid by \obv{multiplicative_monoid_is_exact_monoid}.
    $\carrier[M]=\carrier[R]$ by \auto{multiplicative_monoid_fields}.
    $B\subseteq\carrier[R]$ by \obv{power_bounded_set_carrier}.
    Hence $A\subseteq\carrier[R]$ by \obv{subseteq_transitive}.
    Hence $A,B\subseteq\carrier[M]$.
    $\pointwiseSubsetPowers{M}{A}
        \subseteq\pointwiseSubsetPowers{M}{B}$ by
        \obv{pointwise_subset_powers_monotone}.
    $\Bounded{R}{\pointwiseSubsetPowers{M}{B}}$ by
        \obv{power_bounded_set_powers_bounded}.
    Hence $\Bounded{R}{\pointwiseSubsetPowers{M}{A}}$ by
        \obv{bounded_downward}.
    Follows by \obv{power_bounded_set}.
\end{proof}

\begin{definition}\label{power_bounded_element}
    $x$ is power bounded in $R$ iff
        $R$ is a commutative topological ring and $x\in\carrier[R]$ and
        $\PowerBounded{R}{\{x\}}$.
\end{definition}

\begin{proposition}\label{power_bounded_element_iff}
    Let $R$ be a commutative topological ring.
    Then $x$ is power bounded in $R$ iff
        $x\in\carrier[R]$ and $\PowerBounded{R}{\{x\}}$.
\end{proposition}
\begin{proof}
    Follows by \auto{power_bounded_element}.
\end{proof}

\begin{definition}\label{power_bounded_carrier}
    $\PowerBoundedCarrier{R}=\{x\in\carrier[R]\mid
        \text{$x$ is power bounded in $R$}\}$.
\end{definition}

\begin{proposition}\label{power_bounded_carrier_iff}
    $x\in\PowerBoundedCarrier{R}$ iff
        $x\in\carrier[R]$ and $x$ is power bounded in $R$.
\end{proposition}
\begin{proof}
    Follows by \obv{power_bounded_carrier}.
\end{proof}

\begin{proposition}\label{power_bounded_carrier_subseteq}
    $\PowerBoundedCarrier{R}\subseteq\carrier[R]$.
\end{proposition}
\begin{proof}
    Follows by \obv{power_bounded_carrier_iff,subseteq}.
\end{proof}

\subsection{The elementary power-bounded elements}

\begin{proposition}\label{topological_ring_zero_power_envelope}
    Let $R$ be a commutative topological ring.
    Then $\pointwiseSubsetPowers{\MultiplicativeMonoid{R}}{
        \{\ringZero[R]\}}\subseteq\{\ringZero[R],\ringOne[R]\}$.
\end{proposition}
\begin{proof}
    $R$ is a commutative ring by
        \obv{commutative_topological_ring_is_commutative_ring}.
    Hence $R$ is a ring by \obv{commutative_ring_is_ring}.
    Let $M=\MultiplicativeMonoid{R}$.
    $M$ is a monoid by \obv{multiplicative_monoid_is_exact_monoid}.
    $\carrier[M]=\carrier[R]$ by \auto{multiplicative_monoid_fields}.
    $\ringZero[R]\in\carrier[R]$ by \obv{ring_zero_type}.
    Hence $\ringZero[R]\in\carrier[M]$.
    Show that every element of
        $\pointwiseSubsetPowers{M}{\{\ringZero[R]\}}$ is an element of
        $\{\ringZero[R],\ringOne[R]\}$.
    \begin{subproof}
        Fix $z\in\pointwiseSubsetPowers{M}{\{\ringZero[R]\}}$.
        Take a natural number $n$ such that
            $z\in\pointwiseSubsetPower{M}{\{\ringZero[R]\}}{n}$ by
            \obv{pointwise_subset_powers_iff}.
        $\pointwiseSubsetPower{M}{\{\ringZero[R]\}}{n}
            =\{\monoidPower{M}{\ringZero[R]}{n}\}$ by
            \obv{pointwise_subset_power_singleton}.
        Hence $z=\monoidPower{M}{\ringZero[R]}{n}$ by
            \obv{singleton_iff}.
        $\ringPower{R}{\ringZero[R]}{n}
            =\monoidPower{M}{\ringZero[R]}{n}$ by \obv{ring_power}.
        Hence $z=\ringPower{R}{\ringZero[R]}{n}$.
        \begin{byCase}
        \caseOf{$n=\emptyset$.}
            $\ringPower{R}{\ringZero[R]}{n}=\ringOne[R]$ by
                \auto{ring_power_zero}.
            Hence $\ringOne[R]=z$.
            Follows by \auto{upair_iff}.
        \caseOf{$n\neq\emptyset$.}
            $\ringPower{R}{\ringZero[R]}{n}=\ringZero[R]$ by
                \obv{ring_power_of_zero_positive}.
            Hence $\ringZero[R]=z$.
            Follows by \auto{upair_iff}.
        \end{byCase}
    \end{subproof}
    Hence $\pointwiseSubsetPowers{M}{\{\ringZero[R]\}}
        \subseteq\{\ringZero[R],\ringOne[R]\}$ by \obv{subseteq}.
\end{proof}

\begin{proposition}\label{topological_ring_one_power_envelope}
    Let $R$ be a commutative topological ring.
    Then $\pointwiseSubsetPowers{\MultiplicativeMonoid{R}}{
        \{\ringOne[R]\}}\subseteq\{\ringOne[R]\}$.
\end{proposition}
\begin{proof}
    $R$ is a commutative ring by
        \obv{commutative_topological_ring_is_commutative_ring}.
    Hence $R$ is a ring by \obv{commutative_ring_is_ring}.
    Let $M=\MultiplicativeMonoid{R}$.
    $M$ is a monoid by \obv{multiplicative_monoid_is_exact_monoid}.
    $\carrier[M]=\carrier[R]$ by \auto{multiplicative_monoid_fields}.
    $\ringOne[R]\in\carrier[R]$ by \obv{ring_one_type}.
    Hence $\ringOne[R]\in\carrier[M]$.
    Show that every element of
        $\pointwiseSubsetPowers{M}{\{\ringOne[R]\}}$ is an element of
        $\{\ringOne[R]\}$.
    \begin{subproof}
        Fix $z\in\pointwiseSubsetPowers{M}{\{\ringOne[R]\}}$.
        Take a natural number $n$ such that
            $z\in\pointwiseSubsetPower{M}{\{\ringOne[R]\}}{n}$ by
            \obv{pointwise_subset_powers_iff}.
        $\pointwiseSubsetPower{M}{\{\ringOne[R]\}}{n}
            =\{\monoidPower{M}{\ringOne[R]}{n}\}$ by
            \obv{pointwise_subset_power_singleton}.
        Hence $z=\monoidPower{M}{\ringOne[R]}{n}$ by
            \obv{singleton_iff}.
        $\ringPower{R}{\ringOne[R]}{n}
            =\monoidPower{M}{\ringOne[R]}{n}$ by \obv{ring_power}.
        $\ringPower{R}{\ringOne[R]}{n}=\ringOne[R]$ by
            \obv{ring_power_of_one}.
        Hence $z=\ringOne[R]$.
        Follows by \obv{singleton_iff}.
    \end{subproof}
    Hence $\pointwiseSubsetPowers{M}{\{\ringOne[R]\}}
        \subseteq\{\ringOne[R]\}$ by \obv{subseteq}.
\end{proof}

\begin{theorem}\label{ring_zero_power_bounded}
    Let $R$ be a commutative topological ring.
    Then $\ringZero[R]$ is power bounded in $R$.
\end{theorem}
\begin{proof}
    $\ringZero[R],\ringOne[R]\in\carrier[R]$ by
        \obv{ring_zero_type,ring_one_type}.
    $\{\ringOne[R]\}$ is finite by \obv{singleton_is_finite}.
    Hence $\cons{\ringZero[R]}{\{\ringOne[R]\}}$ is finite by
        \obv{finite_cons}.
    $\{\ringZero[R],\ringOne[R]\}
        =\cons{\ringZero[R]}{\{\ringOne[R]\}}$ by
        \obv{setext,upair_iff,cons_iff,singleton_iff}.
    Hence $\{\ringZero[R],\ringOne[R]\}$ is finite.
    $\{\ringZero[R],\ringOne[R]\}\subseteq\carrier[R]$ by
        \auto{subseteq,upair_iff}.
    Hence $\Bounded{R}{\{\ringZero[R],\ringOne[R]\}}$ by
        \obv{finite_subset_bounded}.
    $\pointwiseSubsetPowers{\MultiplicativeMonoid{R}}{
        \{\ringZero[R]\}}\subseteq\{\ringZero[R],\ringOne[R]\}$ by
        \obv{topological_ring_zero_power_envelope}.
    Hence $\Bounded{R}{\pointwiseSubsetPowers{
        \MultiplicativeMonoid{R}}{\{\ringZero[R]\}}}$ by
        \obv{bounded_downward}.
    $\{\ringZero[R]\}\subseteq\carrier[R]$ by
        \auto{subseteq,singleton_iff}.
    Hence $\PowerBounded{R}{\{\ringZero[R]\}}$ by
        \obv{power_bounded_set}.
    Follows by \obv{power_bounded_element}.
\end{proof}

\begin{theorem}\label{ring_one_power_bounded}
    Let $R$ be a commutative topological ring.
    Then $\ringOne[R]$ is power bounded in $R$.
\end{theorem}
\begin{proof}
    $\ringOne[R]\in\carrier[R]$ by \obv{ring_one_type}.
    $\{\ringOne[R]\}$ is finite by \obv{singleton_is_finite}.
    $\{\ringOne[R]\}\subseteq\carrier[R]$ by
        \auto{subseteq,singleton_iff}.
    Hence $\Bounded{R}{\{\ringOne[R]\}}$ by
        \obv{finite_subset_bounded}.
    $\pointwiseSubsetPowers{\MultiplicativeMonoid{R}}{
        \{\ringOne[R]\}}\subseteq\{\ringOne[R]\}$ by
        \obv{topological_ring_one_power_envelope}.
    Hence $\Bounded{R}{\pointwiseSubsetPowers{
        \MultiplicativeMonoid{R}}{\{\ringOne[R]\}}}$ by
        \obv{bounded_downward}.
    Hence $\PowerBounded{R}{\{\ringOne[R]\}}$ by
        \obv{power_bounded_set}.
    Follows by \obv{power_bounded_element}.
\end{proof}

\begin{proposition}\label{topological_ring_negative_one_power_envelope}
    Let $R$ be a commutative topological ring.
    Then $\pointwiseSubsetPowers{\MultiplicativeMonoid{R}}{
        \{\ringNeg[R](\ringOne[R])\}}
        \subseteq\{\ringOne[R],\ringNeg[R](\ringOne[R])\}$.
\end{proposition}
\begin{proof}
    $R$ is a commutative ring by
        \obv{commutative_topological_ring_is_commutative_ring}.
    Hence $R$ is a ring by \obv{commutative_ring_is_ring}.
    Let $M=\MultiplicativeMonoid{R}$.
    $M$ is a monoid by \obv{multiplicative_monoid_is_exact_monoid}.
    $\carrier[M]=\carrier[R]$ by \auto{multiplicative_monoid_fields}.
    $\ringOne[R]\in\carrier[R]$ by \obv{ring_one_type}.
    Hence $\ringNeg[R](\ringOne[R])\in\carrier[R]$ by \obv{ring_neg_type}.
    Hence $\ringNeg[R](\ringOne[R])\in\carrier[M]$.
    Show that every element of
        $\pointwiseSubsetPowers{M}{\{\ringNeg[R](\ringOne[R])\}}$ is an
        element of $\{\ringOne[R],\ringNeg[R](\ringOne[R])\}$.
    \begin{subproof}
        Fix $z\in
            \pointwiseSubsetPowers{M}{\{\ringNeg[R](\ringOne[R])\}}$.
        Take a natural number $n$ such that
            $z\in\pointwiseSubsetPower{M}{
                \{\ringNeg[R](\ringOne[R])\}}{n}$ by
            \obv{pointwise_subset_powers_iff}.
        $\pointwiseSubsetPower{M}{\{\ringNeg[R](\ringOne[R])\}}{n}
            =\{\monoidPower{M}{\ringNeg[R](\ringOne[R])}{n}\}$ by
            \obv{pointwise_subset_power_singleton}.
        Hence $z=\monoidPower{M}{\ringNeg[R](\ringOne[R])}{n}$ by
            \auto{singleton_iff}.
        $\ringPower{R}{\ringNeg[R](\ringOne[R])}{n}
            =\monoidPower{M}{\ringNeg[R](\ringOne[R])}{n}$ by
            \obv{ring_power}.
        Hence $z=\ringPower{R}{\ringNeg[R](\ringOne[R])}{n}$.
        $\ringPower{R}{\ringNeg[R](\ringOne[R])}{n}
            \in\{\ringOne[R],\ringNeg[R](\ringOne[R])\}$ by
            \obv{ring_power_of_negative_one_in_pair}.
    \end{subproof}
    Hence $\pointwiseSubsetPowers{M}{\{\ringNeg[R](\ringOne[R])\}}
        \subseteq\{\ringOne[R],\ringNeg[R](\ringOne[R])\}$ by
        \obv{subseteq}.
\end{proof}

\begin{theorem}\label{ring_negative_one_power_bounded}
    Let $R$ be a commutative topological ring.
    Then $\ringNeg[R](\ringOne[R])$ is power bounded in $R$.
\end{theorem}
\begin{proof}
    $R$ is a commutative ring by
        \obv{commutative_topological_ring_is_commutative_ring}.
    Hence $R$ is a ring by \obv{commutative_ring_is_ring}.
    $\ringOne[R]\in\carrier[R]$ by \obv{ring_one_type}.
    Hence $\ringNeg[R](\ringOne[R])\in\carrier[R]$ by \obv{ring_neg_type}.
    $\{\ringNeg[R](\ringOne[R])\}$ is finite by
        \obv{singleton_is_finite}.
    Hence $\cons{\ringOne[R]}{\{\ringNeg[R](\ringOne[R])\}}$ is finite by
        \obv{finite_cons}.
    $\{\ringOne[R],\ringNeg[R](\ringOne[R])\}
        =\cons{\ringOne[R]}{\{\ringNeg[R](\ringOne[R])\}}$ by
        \obv{setext,upair_iff,cons_iff,singleton_iff}.
    Hence $\{\ringOne[R],\ringNeg[R](\ringOne[R])\}$ is finite.
    $\{\ringOne[R],\ringNeg[R](\ringOne[R])\}\subseteq\carrier[R]$ by
        \auto{subseteq,upair_iff}.
    Hence $\Bounded{R}{\{\ringOne[R],\ringNeg[R](\ringOne[R])\}}$ by
        \obv{finite_subset_bounded}.
    $\pointwiseSubsetPowers{\MultiplicativeMonoid{R}}{
        \{\ringNeg[R](\ringOne[R])\}}
        \subseteq\{\ringOne[R],\ringNeg[R](\ringOne[R])\}$ by
        \obv{topological_ring_negative_one_power_envelope}.
    Hence $\Bounded{R}{\pointwiseSubsetPowers{
        \MultiplicativeMonoid{R}}{\{\ringNeg[R](\ringOne[R])\}}}$ by
        \obv{bounded_downward}.
    $\{\ringNeg[R](\ringOne[R])\}\subseteq\carrier[R]$ by
        \auto{subseteq,singleton_iff}.
    Hence $\PowerBounded{R}{\{\ringNeg[R](\ringOne[R])\}}$ by
        \obv{power_bounded_set}.
    Follows by \obv{power_bounded_element}.
\end{proof}

\subsection{Multiplication}

\begin{theorem}\label{power_bounded_element_multiplication}
    Suppose $x$ and $y$ are power bounded in $R$.
    Then $\ringMul[R](x,y)$ is power bounded in $R$.
\end{theorem}
\begin{proof}
    $R$ is a commutative topological ring and $x,y\in\carrier[R]$ by
        \obv{power_bounded_element}.
    Hence $R$ is a commutative ring by
        \obv{commutative_topological_ring_is_commutative_ring}.
    Hence $R$ is a ring by \obv{commutative_ring_is_ring}.
    Let $M=\MultiplicativeMonoid{R}$.
    $M$ is a monoid by \obv{multiplicative_monoid_is_exact_monoid}.
    $\carrier[M]=\carrier[R]$ and $\mul[M]=\ringMul[R]$ by
        \auto{multiplicative_monoid_fields}.
    Hence $x,y\in\carrier[M]$.
    $\ringMul[R](x,y)\in\carrier[R]$ by \obv{ring_mul_type}.
    Hence $\ringMul[R](x,y)\in\carrier[M]$.
    $\PowerBounded{R}{\{x\}}$ by \obv{power_bounded_element}.
    Hence $\Bounded{R}{\pointwiseSubsetPowers{
        \MultiplicativeMonoid{R}}{\{x\}}}$ by
        \obv{power_bounded_set_powers_bounded}.
    Hence $\Bounded{R}{\pointwiseSubsetPowers{M}{\{x\}}}$.
    Hence $\pointwiseSubsetPowers{M}{\{x\}}$ is left bounded in $R$ by
        \obv{bounded_in_commutative_topological_ring}.
    $\PowerBounded{R}{\{y\}}$ by \obv{power_bounded_element}.
    Hence $\Bounded{R}{\pointwiseSubsetPowers{
        \MultiplicativeMonoid{R}}{\{y\}}}$ by
        \obv{power_bounded_set_powers_bounded}.
    Hence $\Bounded{R}{\pointwiseSubsetPowers{M}{\{y\}}}$.
    Hence $\pointwiseSubsetPowers{M}{\{y\}}$ is left bounded in $R$ by
        \obv{bounded_in_commutative_topological_ring}.
    Hence $\pointwiseProduct{
        \MultiplicativeMonoid{R}}{
        \pointwiseSubsetPowers{M}{\{x\}}}{
        \pointwiseSubsetPowers{M}{\{y\}}}$ is left bounded in $R$ by
        \obv{left_bounded_pointwise_product}.
    Hence $\pointwiseProduct{M}{
        \pointwiseSubsetPowers{M}{\{x\}}}{
        \pointwiseSubsetPowers{M}{\{y\}}}$ is left bounded in $R$.
    $\mul[M](x,y)=\mul[M](y,x)$ by
        \auto{commutative_ring_mul_commutative}.
    $\pointwiseSubsetPowers{M}{\{\mul[M](x,y)\}}
        \subseteq\pointwiseProduct{M}{
            \pointwiseSubsetPowers{M}{\{x\}}}{
            \pointwiseSubsetPowers{M}{\{y\}}}$ by
        \obv{pointwise_subset_powers_singleton_product}.
    Hence $\pointwiseSubsetPowers{M}{\{\mul[M](x,y)\}}$
        is left bounded in $R$ by \obv{left_bounded_downward}.
    Hence $\Bounded{R}{\pointwiseSubsetPowers{M}{
        \{\mul[M](x,y)\}}}$ by
        \obv{bounded_in_commutative_topological_ring}.
    $\pointwiseSubsetPowers{M}{\{\mul[M](x,y)\}}
        =\pointwiseSubsetPowers{M}{\{\ringMul[R](x,y)\}}$.
    Hence $\Bounded{R}{\pointwiseSubsetPowers{M}{
        \{\ringMul[R](x,y)\}}}$.
    $\{\ringMul[R](x,y)\}\subseteq\carrier[R]$ by
        \auto{subseteq,singleton_iff}.
    Hence $\PowerBounded{R}{\{\ringMul[R](x,y)\}}$ by
        \obv{power_bounded_set}.
    Follows by \obv{power_bounded_element}.
\end{proof}

\subsection{Negation}

\begin{theorem}\label{power_bounded_element_negation}
    Suppose $x$ is power bounded in $R$.
    Then $\ringNeg[R](x)$ is power bounded in $R$.
\end{theorem}
\begin{proof}
    $R$ is a commutative topological ring and $x\in\carrier[R]$ by
        \obv{power_bounded_element}.
    Hence $R$ is a commutative ring by
        \obv{commutative_topological_ring_is_commutative_ring}.
    Hence $R$ is a ring by \obv{commutative_ring_is_ring}.
    $\ringOne[R]\in\carrier[R]$ by \obv{ring_one_type}.
    Hence $\ringNeg[R](\ringOne[R])$ is power bounded in $R$ by
        \obv{ring_negative_one_power_bounded}.
    Hence $\ringMul[R](\ringNeg[R](\ringOne[R]),x)$ is power bounded in $R$
        by \obv{power_bounded_element_multiplication}.
    $\ringMul[R](\ringNeg[R](\ringOne[R]),x)
        =\ringNeg[R](\ringMul[R](\ringOne[R],x))$ by
        \obv{ring_neg_product_left}.
    $\ringMul[R](\ringOne[R],x)=x$ by \obv{ring_mul_one_left}.
    Hence $\ringMul[R](\ringNeg[R](\ringOne[R]),x)
        =\ringNeg[R](x)$.
\end{proof}

\subsection{Addition}

\begin{theorem}\label{power_bounded_element_addition}
    Let $R$ be a nonarchimedean topological ring.
    Suppose $x$ and $y$ are power bounded in $R$.
    Then $\ringAdd[R](x,y)$ is power bounded in $R$.
\end{theorem}
\begin{proof}
    $R$ is a commutative topological ring and $x,y\in\carrier[R]$ by
        \obv{power_bounded_element}.
    Hence $R$ is a commutative ring by
        \obv{commutative_topological_ring_is_commutative_ring}.
    Hence $R$ is a ring by \obv{commutative_ring_is_ring}.
    Let $M=\MultiplicativeMonoid{R}$.
    $M$ is a monoid by \obv{multiplicative_monoid_is_exact_monoid}.
    $\carrier[M]=\carrier[R]$ by \auto{multiplicative_monoid_fields}.
    Hence $x,y\in\carrier[M]$.
    $\ringAdd[R](x,y)\in\carrier[R]$ by \obv{ring_add_type}.
    Hence $\ringAdd[R](x,y)\in\carrier[M]$.
    $\PowerBounded{R}{\{x\}}$ by \obv{power_bounded_element}.
    Hence $\Bounded{R}{\pointwiseSubsetPowers{
        \MultiplicativeMonoid{R}}{\{x\}}}$ by
        \obv{power_bounded_set_powers_bounded}.
    Hence $\Bounded{R}{\pointwiseSubsetPowers{M}{\{x\}}}$.
    Hence $\pointwiseSubsetPowers{M}{\{x\}}$ is left bounded in $R$ by
        \obv{bounded_in_commutative_topological_ring}.
    $\PowerBounded{R}{\{y\}}$ by \obv{power_bounded_element}.
    Hence $\Bounded{R}{\pointwiseSubsetPowers{
        \MultiplicativeMonoid{R}}{\{y\}}}$ by
        \obv{power_bounded_set_powers_bounded}.
    Hence $\Bounded{R}{\pointwiseSubsetPowers{M}{\{y\}}}$.
    Hence $\pointwiseSubsetPowers{M}{\{y\}}$ is left bounded in $R$ by
        \obv{bounded_in_commutative_topological_ring}.
    Let $X=\pointwiseSubsetPowers{M}{\{x\}}$ and
        $Y=\pointwiseSubsetPowers{M}{\{y\}}$.
    Hence $X$ is left bounded in $R$.
    Hence $Y$ is left bounded in $R$.
    Hence $\pointwiseProduct{\MultiplicativeMonoid{R}}{X}{Y}$ is left
        bounded in $R$ by \obv{left_bounded_pointwise_product}.
    Hence $\generatedSubgroup{
        \pointwiseProduct{\MultiplicativeMonoid{R}}{X}{Y}}{
        \AdditiveGroup{R}}$ is left bounded in $R$ by
        \obv{nonarchimedean_generated_subgroup_left_bounded}.
    $\additiveSubsetProduct{R}{X}{Y}
        =\generatedSubgroup{
            \pointwiseProduct{\MultiplicativeMonoid{R}}{X}{Y}}{
            \AdditiveGroup{R}}$ by \obv{additive_subset_product}.
    Hence $\additiveSubsetProduct{R}{X}{Y}$ is left bounded in $R$.
    Show that $\pointwiseSubsetPowers{M}{\{\ringAdd[R](x,y)\}}
        \subseteq\additiveSubsetProduct{R}{X}{Y}$.
    \begin{subproof}
        Show that every element of
            $\pointwiseSubsetPowers{M}{\{\ringAdd[R](x,y)\}}$ is an
            element of $\additiveSubsetProduct{R}{X}{Y}$.
        \begin{subproof}
            Fix $z\in\pointwiseSubsetPowers{M}{\{\ringAdd[R](x,y)\}}$.
            Take a natural number $n$ such that
                $z\in\pointwiseSubsetPower{M}{
                    \{\ringAdd[R](x,y)\}}{n}$ by
                \obv{pointwise_subset_powers_iff}.
            $\pointwiseSubsetPower{M}{\{\ringAdd[R](x,y)\}}{n}
                =\{\monoidPower{M}{\ringAdd[R](x,y)}{n}\}$ by
                \obv{pointwise_subset_power_singleton}.
            Hence $z=\monoidPower{M}{\ringAdd[R](x,y)}{n}$ by
                \auto{singleton_iff}.
            $\ringPower{R}{\ringAdd[R](x,y)}{n}
                =\monoidPower{M}{\ringAdd[R](x,y)}{n}$ by
                \obv{ring_power}.
            Hence $z=\ringPower{R}{\ringAdd[R](x,y)}{n}$.
            $\ringPower{R}{\ringAdd[R](x,y)}{n}\in
                \additiveSubsetProduct{R}{
                    \pointwiseSubsetPowers{\MultiplicativeMonoid{R}}{
                        \{x\}}}{
                    \pointwiseSubsetPowers{\MultiplicativeMonoid{R}}{
                        \{y\}}}$ by
                \obv{ring_power_of_sum_in_power_envelope_product}.
            Hence $\ringPower{R}{\ringAdd[R](x,y)}{n}
                \in\additiveSubsetProduct{R}{X}{Y}$.
        \end{subproof}
        Follows by \obv{subseteq}.
    \end{subproof}
    Hence $\pointwiseSubsetPowers{\MultiplicativeMonoid{R}}{
        \{\ringAdd[R](x,y)\}}
        \subseteq\additiveSubsetProduct{R}{X}{Y}$.
    Hence $\pointwiseSubsetPowers{\MultiplicativeMonoid{R}}{
        \{\ringAdd[R](x,y)\}}$
        is left bounded in $R$ by \obv{left_bounded_downward}.
    Hence $\Bounded{R}{\pointwiseSubsetPowers{\MultiplicativeMonoid{R}}{
        \{\ringAdd[R](x,y)\}}}$ by
        \obv{bounded_in_commutative_topological_ring}.
    $\{\ringAdd[R](x,y)\}\subseteq\carrier[R]$ by
        \auto{subseteq,singleton_iff}.
    Hence $\PowerBounded{R}{\{\ringAdd[R](x,y)\}}$ by
        \obv{power_bounded_set}.
    Follows by \obv{power_bounded_element}.
\end{proof}

\section{The ring of power-bounded elements}

\begin{proposition}\label{ring_zero_in_power_bounded_carrier}
    Let $R$ be a commutative topological ring.
    Then $\ringZero[R]\in\PowerBoundedCarrier{R}$.
\end{proposition}
\begin{proof}
    $\ringZero[R]\in\carrier[R]$ by
        \obv{commutative_topological_ring_is_commutative_ring,
            commutative_ring_is_ring,ring_zero_type}.
    $\ringZero[R]$ is power bounded in $R$ by
        \obv{ring_zero_power_bounded}.
    Follows by \obv{power_bounded_carrier_iff}.
\end{proof}

\begin{proposition}\label{ring_one_in_power_bounded_carrier}
    Let $R$ be a commutative topological ring.
    Then $\ringOne[R]\in\PowerBoundedCarrier{R}$.
\end{proposition}
\begin{proof}
    $\ringOne[R]\in\carrier[R]$ by
        \obv{commutative_topological_ring_is_commutative_ring,
            commutative_ring_is_ring,ring_one_type}.
    $\ringOne[R]$ is power bounded in $R$ by
        \obv{ring_one_power_bounded}.
    Follows by \obv{power_bounded_carrier_iff}.
\end{proof}

\begin{proposition}\label{power_bounded_carrier_negation}
    Suppose $x\in\PowerBoundedCarrier{R}$.
    Then $\ringNeg[R](x)\in\PowerBoundedCarrier{R}$.
\end{proposition}
\begin{proof}
    $x\in\carrier[R]$ and $x$ is power bounded in $R$ by
        \obv{power_bounded_carrier_iff}.
    Hence $R$ is a commutative ring by
        \obv{commutative_topological_ring_is_commutative_ring}.
    Hence $R$ is a ring by \obv{commutative_ring_is_ring}.
    $\ringNeg[R](x)\in\carrier[R]$ by \obv{ring_neg_type}.
    $\ringNeg[R](x)$ is power bounded in $R$ by
        \obv{power_bounded_element_negation}.
    Follows by \obv{power_bounded_carrier_iff}.
\end{proof}

\begin{proposition}\label{power_bounded_carrier_multiplication}
    Suppose $x,y\in\PowerBoundedCarrier{R}$.
    Then $\ringMul[R](x,y)\in\PowerBoundedCarrier{R}$.
\end{proposition}
\begin{proof}
    $x,y\in\carrier[R]$ and $x$ and $y$ are power bounded in $R$ by
        \obv{power_bounded_carrier_iff}.
    Hence $R$ is a commutative ring by
        \obv{commutative_topological_ring_is_commutative_ring}.
    Hence $R$ is a ring by \obv{commutative_ring_is_ring}.
    $\ringMul[R](x,y)\in\carrier[R]$ by \obv{ring_mul_type}.
    $\ringMul[R](x,y)$ is power bounded in $R$ by
        \obv{power_bounded_element_multiplication}.
    Follows by \obv{power_bounded_carrier_iff}.
\end{proof}

\begin{proposition}\label{power_bounded_carrier_addition}
    Let $R$ be a nonarchimedean topological ring.
    Suppose $x,y\in\PowerBoundedCarrier{R}$.
    Then $\ringAdd[R](x,y)\in\PowerBoundedCarrier{R}$.
\end{proposition}
\begin{proof}
    $x,y\in\carrier[R]$ and $x$ and $y$ are power bounded in $R$ by
        \obv{power_bounded_carrier_iff}.
    Hence $R$ is a commutative ring by
        \obv{commutative_topological_ring_is_commutative_ring}.
    Hence $R$ is a ring by \obv{commutative_ring_is_ring}.
    $\ringAdd[R](x,y)\in\carrier[R]$ by \obv{ring_add_type}.
    $\ringAdd[R](x,y)$ is power bounded in $R$ by
        \obv{power_bounded_element_addition}.
    Follows by \obv{power_bounded_carrier_iff}.
\end{proof}

\begin{theorem}\label{power_bounded_carrier_is_subring}
    Let $R$ be a commutative topological ring.
    Suppose $R$ is a nonarchimedean topological ring.
    Then $\PowerBoundedCarrier{R}$ is a subring of $R$.
\end{theorem}
\begin{proof}
    $R$ is a commutative ring by
        \obv{commutative_topological_ring_is_commutative_ring}.
    Hence $R$ is a ring by \obv{commutative_ring_is_ring}.
    Let $G=\AdditiveGroup{R}$.
    $G$ is a group by \obv{additive_group_is_group}.
    $\carrier[G]=\carrier[R]$ and $\mul[G]=\ringAdd[R]$ and
        $\neutral[G]=\ringZero[R]$ and $\groupInv[G]=\ringNeg[R]$ by
        \auto{additive_group_fields}.
    $\PowerBoundedCarrier{R}\subseteq\carrier[R]$ by
        \obv{power_bounded_carrier_subseteq}.
    Hence $\PowerBoundedCarrier{R}\subseteq\carrier[G]$.
    $\ringZero[R]\in\PowerBoundedCarrier{R}$ by
        \obv{ring_zero_in_power_bounded_carrier}.
    Hence $\neutral[G]\in\PowerBoundedCarrier{R}$.
    For all $x,y\in\PowerBoundedCarrier{R}$ we have
        $\ringAdd[R](x,y)\in\PowerBoundedCarrier{R}$ by
        \obv{power_bounded_carrier_addition}.
    Hence for all $x,y\in\PowerBoundedCarrier{R}$ we have
        $\mul[G](x,y)\in\PowerBoundedCarrier{R}$ by \auto.
    For every $x\in\PowerBoundedCarrier{R}$ we have
        $\ringNeg[R](x)\in\PowerBoundedCarrier{R}$ by
        \obv{power_bounded_carrier_negation}.
    Hence for every $x\in\PowerBoundedCarrier{R}$ we have
        $\groupInv[G](x)\in\PowerBoundedCarrier{R}$ by \auto.
    Hence $\subgroup{\PowerBoundedCarrier{R}}{G}$ by \obv{subgroup}.
    Hence $\subgroup{\PowerBoundedCarrier{R}}{\AdditiveGroup{R}}$.
    $\ringOne[R]\in\PowerBoundedCarrier{R}$ by
        \obv{ring_one_in_power_bounded_carrier}.
    For all $x,y\in\PowerBoundedCarrier{R}$ we have
        $\ringMul[R](x,y)\in\PowerBoundedCarrier{R}$ by
        \obv{power_bounded_carrier_multiplication}.
    Follows by \obv{subring}.
\end{proof}

\begin{definition}\label{power_bounded_ring}
    $\PowerBoundedRing{R}
        =\InducedSubring{R}{\PowerBoundedCarrier{R}}$.
\end{definition}

\begin{theorem}\label{power_bounded_ring_agreement}
    Let $R$ be a commutative topological ring.
    Suppose $R$ is a nonarchimedean topological ring.
    Then $\carrier[\PowerBoundedRing{R}]=\PowerBoundedCarrier{R}$ and
        $\ringZero[\PowerBoundedRing{R}]=\ringZero[R]$ and
        $\ringOne[\PowerBoundedRing{R}]=\ringOne[R]$ and
        (for all $x,y\in\PowerBoundedCarrier{R}$ we have
            $\ringAdd[\PowerBoundedRing{R}](x,y)=\ringAdd[R](x,y)$ and
            $\ringMul[\PowerBoundedRing{R}](x,y)=\ringMul[R](x,y)$) and
        (for every $x\in\PowerBoundedCarrier{R}$ we have
            $\ringNeg[\PowerBoundedRing{R}](x)=\ringNeg[R](x)$).
\end{theorem}
\begin{proof}
    $\PowerBoundedCarrier{R}$ is a subring of $R$ by
        \obv{power_bounded_carrier_is_subring}.
    Follows by \auto{power_bounded_ring,induced_subring_agreement}.
\end{proof}

\begin{theorem}\label{power_bounded_ring_is_exact_commutative_ring}
    Let $R$ be a commutative topological ring.
    Suppose $R$ is a nonarchimedean topological ring.
    Then $\PowerBoundedRing{R}$ is an exact ring and
        $\PowerBoundedRing{R}$ is a commutative ring.
\end{theorem}
\begin{proof}
    Let $S=\PowerBoundedRing{R}$.
    $\PowerBoundedCarrier{R}$ is a subring of $R$ by
        \obv{power_bounded_carrier_is_subring}.
    Hence $\InducedSubring{R}{\PowerBoundedCarrier{R}}$ is an exact ring by
        \obv{induced_subring_is_exact}.
    Hence $S$ is an exact ring by \auto{power_bounded_ring}.
    Show that for all $x,y\in\carrier[S]$ we have
        $\ringMul[S](x,y)=\ringMul[S](y,x)$.
    \begin{subproof}
        Fix $x,y\in\carrier[S]$.
        $\carrier[S]=\PowerBoundedCarrier{R}$ by
            \auto{power_bounded_ring_agreement}.
        Hence $x,y\in\PowerBoundedCarrier{R}$.
        $\ringMul[S](x,y)=\ringMul[R](x,y)$ and
            $\ringMul[S](y,x)=\ringMul[R](y,x)$ by
            \auto{power_bounded_ring_agreement}.
        $x,y\in\carrier[R]$ by
            \auto{power_bounded_carrier_iff}.
        Follows by \obv{commutative_ring_mul_commutative}.
    \end{subproof}
    Hence $S$ is a commutative ring by \obv{commutative_ring_intro}.
    Follows by \auto.
\end{proof}

\begin{classification}\label{power_bounded_ring_is_commutative_ring}
    Let $R$ be a commutative topological ring.
    Suppose $R$ is a nonarchimedean topological ring.
    Then $\PowerBoundedRing{R}$ is a commutative ring.
\end{classification}
\begin{proof}
    Follows by \obv{power_bounded_ring_is_exact_commutative_ring}.
\end{proof}

\begin{theorem}\label{power_bounded_ring_natural_agreement}
    Let $R$ be a commutative topological ring.
    Suppose $R$ is a nonarchimedean topological ring.
    Let $n$ be a natural number.
    Then $\ringNatural{\PowerBoundedRing{R}}{n}=\ringNatural{R}{n}$.
\end{theorem}
\begin{proof}
    $\PowerBoundedCarrier{R}$ is a subring of $R$ by
        \obv{power_bounded_carrier_is_subring}.
    $\ringNatural{\InducedSubring{R}{\PowerBoundedCarrier{R}}}{n}
        =\ringNatural{R}{n}$ by
        \obv{induced_subring_natural_agreement}.
    Follows by \auto{power_bounded_ring}.
\end{proof}

\begin{theorem}\label{power_bounded_ring_power_agreement}
    Let $R$ be a commutative topological ring.
    Suppose $R$ is a nonarchimedean topological ring.
    Suppose $x\in\PowerBoundedCarrier{R}$.
    Let $n$ be a natural number.
    Then $\ringPower{\PowerBoundedRing{R}}{x}{n}
        =\ringPower{R}{x}{n}$.
\end{theorem}
\begin{proof}
    $\PowerBoundedCarrier{R}$ is a subring of $R$ by
        \obv{power_bounded_carrier_is_subring}.
    $\ringPower{\InducedSubring{R}{\PowerBoundedCarrier{R}}}{x}{n}
        =\ringPower{R}{x}{n}$ by
        \obv{induced_subring_power_agreement}.
    Follows by \auto{power_bounded_ring}.
\end{proof}

\begin{definition}\label{uniform_topological_ring}
    $R$ is a uniform topological ring iff
        $R$ is a commutative topological ring and
        $\Bounded{R}{\PowerBoundedCarrier{R}}$.
\end{definition}
